\documentclass[a4paper,12pt]{article}

\usepackage[utf8]{inputenc}
\usepackage[T2A]{fontenc}
\usepackage[russian]{babel}
\usepackage{geometry}
\usepackage{array}
\usepackage{longtable}

\geometry{margin=2cm}

\title{Портативный тепличный комплекс\\
с полным контролем микроклимата}
\author{
Ментор: Павел Владимирович \\
Команда:    Загороднюк Михаил, Б01-503, @mister\_mikl \\
            Чепайкин Семен, Б01-503, @Chipulikk \\
            Скворцов Николай, Б01-503, @n1sk\_0v
}
\date{}

\begin{document}

\maketitle

\section{Аннотация}

Проект направлен на разработку портативного тепличного комплекса, обеспечивающего автоматизированный контроль и регулирование агротехнических параметров микроклимата системы: влажности, температуры, освещённости и влажности почвы.

\section{Цель проекта}

Разработка портативного тепличного комплекса с замкнутым контуром управления микроклиматом, обеспечивающего следующие возможности:

\begin{itemize}
    \item автоматическое поддержание заданной влажности и температуры воздуха;
    \item циклическое увлажнение почвы до необходимого уровня;
    \item минимальное участие пользователя в работе системы;
    \item настройка параметров комплекса для различных агротехнических условий растений через веб-интерфейс.
\end{itemize}

\section{Технические требования}

\subsection{Диапазоны регулирования}

\begin{longtable}{|p{6cm}|p{8cm}|}
\hline
\textbf{Параметр} & \textbf{Диапазон} \\
\hline
Влажность воздуха & от комнатной (~30--40\%) до 95\% \\
\hline
Температура воздуха & +15\ldots+40~$^\circ$C \\
\hline
Освещённость & регулируемая (искусственное освещение) \\
\hline
Влажность почвы & автоматическое поддержание \\
\hline
Время автономной работы & $\geq 72$ часа (без вмешательства пользователя) \\
\hline
\end{longtable}

\subsection{Габариты комплекса}

\begin{itemize}
    \item Максимальные размеры установки: 800 $\times$ 800 $\times$ 1000 мм;
    \item Тип конструкции: полугерметичное основание и герметичный внешний корпус;
    \item Материал: фанера или пластик для основания, акрил для корпуса.
\end{itemize}

\section{Функционал системы}

Система должна обеспечивать следующие функции:

\subsection{Контроль параметров среды}

\begin{itemize}
    \item измерение температуры воздуха;
    \item измерение относительной влажности воздуха;
    \item измерение освещённости;
    \item измерение влажности почвы;
    \item контроль уровня воды в резервуаре.
\end{itemize}

\subsection{Регулирование}

\begin{itemize}
    \item управление системой увлажнения воздуха;
    \item управление вентиляцией (приток/вытяжка);
    \item управление нагревом (при необходимости);
    \item управление освещением;
    \item автоматический полив растений.
\end{itemize}

\subsection{Автоматизация}

\begin{itemize}
    \item работа по заданным уставкам;
    \item возможность настройки пользователем;
    \item поддержание параметров по обратной связи (замкнутый контур управления).
\end{itemize}

\section{Задачи проекта}

\begin{itemize}
    \item анализ и выбор элементной базы;
    \item разработка алгоритмов управления микроклиматом;
    \item проектирование электрической схемы;
    \item разработка CAD-модели конструкции;
    \item проведение CAE-расчётов (тепловые потоки, циркуляция воздуха);
    \item изготовление прототипа;
    \item сборка системы;
    \item проведение испытаний и калибровка.
\end{itemize}

\section{Элементная база (предварительная)}

\subsection{Контроллер}

\begin{itemize}
    \item Arduino Mini / Arduino Nano / ESP32 (предпочтительно для расширения функционала).
\end{itemize}

\subsection{Сенсоры}

\begin{itemize}
    \item датчик температуры и влажности воздуха (например: DHT22 / SHT31);
    \item датчик освещённости (LDR / BH1750);
    \item датчик влажности почвы (ёмкостной тип);
    \item датчик уровня воды.
\end{itemize}

\subsection{Исполнительные устройства}

\begin{itemize}
    \item вентиляторы (DC);
    \item насос (перистальтический или погружной);
    \item увлажнитель (ультразвуковой модуль или распылитель);
    \item светодиодные модули (LED grow light).
\end{itemize}

\subsection{Дополнительные компоненты}

\begin{itemize}
    \item воздушный фильтр (для приточного воздуха);
    \item резервуар для воды;
    \item воздуховоды / каналы;
    \item корпус (акрил/поликарбонат).
\end{itemize}

\end{document}
